% v25: unbounded endpoint exclusions on the proportional zero axis.
\section{比例零軸の両端点における無上界排除と固定一整数での有限復元}
\label{hpe25:section}

本節は比例零軸 $x=y=0$ の端点 $n=0,z$ を扱う．
定理~\ref{hp24:main} の必要条件
\[
 s<t,\qquad b\ge2,\qquad z\ge1,\qquad 0\le n\le z
\]
を前提とする．記号を
\[
 F=b-1+sz,\quad m=F^2+F+1,\quad K=s(F-1)-z+1,\quad
 C_n=n((sz)^2-3sz+3)+(s^2-1)z-3s+1
\]
とし，差方程式と元の pair-product 残余を
\begin{align}
 (b+1)t&=mn+K,\label{hpe25:Q}\\
 \mathfrak r&=-bt^2+(s^2z-2s-bz)t+F(s^2-b)
 +n\{sm+(F^2+b)t\}=0\label{hpe25:residual}
\end{align}
と書く．既知の $C_n=0$ の二つの例外は模型を与えないので，
以下の模型では必ず $C_n>0$ かつ $b+1\mid C_n$ である．

\begin{theorem}[比例零軸の端点有限化]\label{hpe25:main}
正の canonical 整数模型に対し，次が必要である．
\begin{enumerate}
\item $n=0$ なら
\[
 s\ge3,\qquad s<b<\frac32s,\qquad 1\le z<\frac s2.
\]
\item $n=z$ なら
\[
 s\ge2,\qquad 1\le z\le s,\qquad
 1\le e:=\frac{C_z}{b+1}\le s-1.
\]
\end{enumerate}
従って両端点は，低行の一整数 $s$ だけを固定すれば全て有限に復元できる．
$n=0$ では各 $b$ について $z$ の二次方程式を解くことで，
完全検査前の $(b,z)$ の候補数は高々
$2\lfloor(s-1)/2\rfloor$ である．
$n=z$ では $(z,e)$ の候補数は高々 $s(s-1)$ である．
いずれも $t,S,T,U$ に事前の上限を置かない．
本定理は全ての $s$ を通じた両端点の不存在を主張せず，
内部 $0<n<z$ の一般不存在も主張しない．
\end{theorem}

\subsection{整数商と元の残余}

正整数 $e=C_n/(b+1)$ を導入すると，式~\eqref{hpe25:Q} は
\begin{equation}\label{hpe25:e-coordinates}
 t=s+n(b+2sz-2)+e,\qquad S=s+n(sz-2)+e
\end{equation}
となる．ここで $T=s+Fn$，$U=s+t+z-1-n$ である．
$e$ は独立に選べる正実数ではなく，$b+1$ による正整数商である．
以下の四次式は式~\eqref{hpe25:residual} を代入して作り，
差方程式だけから候補を模型として認めることはしない．

\subsection{端点 \texorpdfstring{$n=0$}{n=0} の三つの無上界排除}

$n=0$ では $S=t,T=s$ なので，正の matching entry
$p_{ba}=s^2-b>0$ より $s\ge2$ かつ $2\le b\le s^2-1$ である．
$t=K/(b+1)$ を元の残余に代入し，
\[
 P_0(s,b,z):=(b+1)^2\mathfrak r
            =A_2z^2+A_1z+A_0
\]
と置くと
\begin{align}
 A_2={}&(s^2-1)(s^2-b^2),\label{hpe25:P0coeff}\\
 A_1={}&(b-1)(-2b^2s-bs-b+3s^3-s^2-2s),\notag\\
 A_0={}&-b^4-b^3+3b^2s^2-2b^2s+b^2-3bs^2+2bs+3s^2-2s.\notag
\end{align}
$s=2$ では $b=2,3$ のみであり，それぞれ
\[
 P_0(2,2,z)=-2(3z-2),\qquad
 P_0(2,3,z)=-(z+1)(15z+43)
\]
となるので，$z\ge1$ に零点はない．従って $s\ge3$ としてよい．

まず $b\le s$ を排除する．$s=b+a$ と書くと
\begin{align*}
P_0={}&z^2a^4+z(4bz+3b-3)a^3\\
 &+\{(5b^2-1)z^2+(9b^2-10b+1)z+3b^2-3b+3\}a^2\\
 &+\{2b(b^2-1)z^2+(7b^3-10b^2+b+2)z
             +6b^3-8b^2+8b-2\}a\\
 &+b(b-1)\{(b-3)(b+1)z+2(b-1)^2\}.
\end{align*}
$b\ge3$ では全ての項が非負で，定数項は正である．
$b=2$ では $s\ge3$ より $a\ge1$ であり，
一次項だけでも負の定数項 $-6z+4$ を上回る：
\[
 (12z^2+20z+30)a-6z+4\ge12z^2+14z+34>0.
\]
従って $P_0=0$ は $b>s$ を必要とする．

次に $b\ge3s/2$ を排除する．$b=3s/2+a$，$a\ge0$ と置くと
\[
 -16P_0=\sum_{i=0}^4 L_i(s,z)a^i
\]
であり，係数は
\begin{align*}
 L_4={}&16,\qquad L_3=16(2sz+6s+1),\\
 L_2={}&8\{2(s^2-1)z^2+(18s^2-2s+2)z+21s^2+13s-2\},\\
 L_1={}&4\{12s(s^2-1)z^2+(42s^3-8s^2+16s-4)z
                    +18s^3+63s^2-20s\},\\
 L_0={}&s\{20s(s^2-1)(z-1)^2
 +(76s^3+36s^2+4s-56)(z-1)\\
 &\hspace{30mm}+29s^3+234s^2-108s-24\}.
\end{align*}
これらは $s\ge3,z\ge1$ で全て正なので $P_0<0$ となる．

最後に $z\ge s/2$ を排除する．既に $b>s$ なので整数
$a=b-s-1\ge0$ を用い，$w=z-s/2\ge0$ と置く．直接展開により
\begin{equation}\label{hpe25:wide-z}
 -4P_0(s,s+1+a,s/2+w)=J_2w^2+J_1w+J_0
\end{equation}
であり，
\begin{align*}
 J_2={}&4(a+1)(s^2-1)(a+2s+1),\\
 J_1={}&4\{2sa^3+(s^3+6s^2+4s+1)a^2
 +(2s^4+5s^3+9s^2+5s+1)a+s^2(s^2+7s+4)\},\\
 J_0={}&4a^4+4(s^2+4s+5)a^3
 +(s^4+12s^3+21s^2+70s+32)a^2\\
 &+2(s^5+4s^4+6s^3+38s^2+37s+10)a
 +5s^4+26s^3+29s^2+28s+4.
\end{align*}
$s\ge3$ では右辺は正なので，この領域も零点を持たない．
これで第一の必要条件が証明された．さらに $b>s$ より
式~\eqref{hpe25:P0coeff} の $A_2$ は負であり，固定 $(s,b)$ ごとの
$z$ の候補は高々二つである．整数 $b$ の個数は
$\lfloor(s-1)/2\rfloor$ なので，定理の候補数が従う．

\subsection{端点 \texorpdfstring{$n=z$}{n=z} の単調四次式}

$n=z$ では
\begin{align*}
 C&:=C_z=s^2z(z^2+1)-3s(z^2+1)+2z+1,\\
 A&:=3s(z^2+1)-4z-1,\\
 B&:=3s^2(z^2+1)^2-2s(z^2+1)(4z+1)+5z^2+4z-2,\\
 H&:=sz(z^2+1)-z^2-z+3
\end{align*}
と置く．式~\eqref{hpe25:e-coordinates} と $b=C/e-1$ を元の残余へ
代入すると，恒等的に
\begin{equation}\label{hpe25:R}
 R(e):=e^2\mathfrak r=e^4+Ae^3+Be^2+CHe-C^2
\end{equation}
となる．従って $e$ はこのモニック四次式の正整数零点でなければならない．

まず $s\ge2$ とする．$a=s-2\ge0$ と置くと
\[
 A=3a(z^2+1)+6z^2-4z+5>0,
\]
また
\begin{align*}
 B={}&3a^2(z^2+1)^2+2a(z^2+1)(6z^2-4z+5)\\
 &+12(z-1)^4+32(z-1)^3+49(z-1)^2+38(z-1)+15>0.
\end{align*}
さらに $H\ge2z^3-z^2+z+3>0$ である．$C>0$ と合わせて
$R'(e)>0$ が $e\ge0$ の全域で成立する．$R(0)=-C^2<0$ なので，
$R$ は正実零点をちょうど一つ持つ．

この零点が $s$ より小さいことを示す．$a=s-2$ に関して
$R(s)=\sum_{i=0}^4 Q_i(z)a^i$ と展開すると
\begin{align*}
Q_0={}&16z^5+8z^4+28z^3+38z^2+6z+25,\\
Q_1={}&24z^5+52z^4+18z^3+171z^2-31z+117,\\
Q_2={}&12z^5+58z^4+2z^3+177z^2-30z+130,\\
Q_3={}&2z^5+23z^4+69z^2-6z+53,\\
Q_4={}&3z^4+9z^2+7.
\end{align*}
$z\ge1$ で各係数が正であるから $R(s)>0$ である．
従って正整数零点は $1\le e\le s-1$ を満たす．

次に $z\ge s+1$ を排除する．$w=z-s-1\ge0$ とすると
\begin{equation}\label{hpe25:root-gap}
 -R(s-1)=\sum_{i=0}^6 E_i(s)w^i,
\end{equation}
ここで
\begin{align*}
 E_6={}&s^3,\\
 E_5={}&2s^2(3s^2+2s-2),\\
 E_4={}&s(15s^4+17s^3-6s^2-16s+5),\\
 E_3={}&2(10s^6+14s^5+2s^4-14s^3-13s^2+11s-1),\\
 E_2={}&15s^7+22s^6+8s^5-15s^4-34s^3-15s^2+38s-10,\\
 E_1={}&(s-1)(6s^7+14s^6+14s^5+4s^4+4s^3-4s^2-24s+13),\\
 E_0={}&(s-1)^2(s^3+2s^2+s-1)(s^4+s^3-4s+1).
\end{align*}
これらは全て $s\ge2$ で正である．例えば $E_2$ では
$15s^7-15s^4$，$22s^6-34s^3$，$8s^5-15s^2$，$38s-10$
をそれぞれ正の組として評価できる．他の係数も同様に，または
$s=a+2$ に展開して全係数が正であることから従う．
従って $R(s-1)<0<R(s)$ であり，唯一の正実零点は開区間
$(s-1,s)$ にある．正整数零点は存在しない．これで $z\le s$ を得る．

残る $s=1$ では $C=(z-1)^3-1$ なので $C>0$ は $z\ge3$ を要求する．
$z\ge3$ で $A,B,H$ は正であり，$B$ は $w=z-2$ と書けば
$3w^4+16w^3+33w^2+32w+11$ である．また
\[
 R(1)=2w^5+14w^4+44w^3+72w^2+61w+18>0.
\]
従って唯一の正実零点は $(0,1)$ にあり，正整数解はない．
定理第二項と候補数 $s(s-1)$ が従う．

\subsection{完全復元と検証の境界}

$n=0$ では定理の有限範囲内で $b$ を列挙し，
式~\eqref{hpe25:P0coeff} の整数零点 $z$ を判別式の完全平方検査で求める．
その後 $b+1\mid K$，$t=K/(b+1)>s$ と元の残余を検査する．
$n=z$ では定理の $(z,e)$ を列挙し，$e\mid C$，$b=C/e-1\ge2$ として
$t=s+z(b+2sz-2)+e$ を求め，式~\eqref{hpe25:R} と元の残余を検査する．

いずれも通過した候補に対し，全ての $S,T,U$ の正値性，符号，
四つの行和，六つの pair-product，欠損
$(D_a,D_b,D_c,D_0)=(0,0,z,-z)$，正の重みの全有向不等式を
有理数演算で再検査して初めて模型と認める．
同梱の \nolinkurl{verify_high_d_axis_endpoints_v25.py} は標準ライブラリのみを用い，
元の残余から恒等式を独立に再構成し，上記の正係数証明を係数比較する．
これは有界探索による無上界主張ではない．
付随する $1\le s\le400$ の端点復元で模型が零個であったことは，
この新しい定理を使った有限検査結果としてのみ記録する．
一般の $s$，および内部 $0<n<z$ は依然として開いている．
